\chapter{Testing from functions to hostile streams}
\chapterintro{The happy path proves almost nothing about a protocol implementation. Confidence comes from testing boundaries: fragmentation, ambiguity, limits, concurrency, rollback, and shutdown.}

\begin{objectives}
\begin{itemize}
  \item Organize model, parser, application, repository, API, and network tests.
  \item Test stream fragmentation independently from packet timing.
  \item Use temporary databases and ephemeral ports for isolated integration tests.
  \item Harden behavior through explicit malformed-input and concurrency cases.
\end{itemize}
\end{objectives}

\section{A testing pyramid for this project}

Fast unit tests cover deterministic functions and objects: Headers, JSON factories,
parser feeds, router precedence, validation, and repository operations. Application
tests call \code{app.handle} directly with Request objects. Network tests start a
real listener on loopback and use raw sockets or \code{http.client}.

Most cases belong below the network boundary because unit tests can control every
fragment without timing. A smaller set of socket tests confirms that accept,
timeouts, keep-alive, serialization, and shutdown connect correctly.

\section{Parser tests are protocol examples}

The parser suite demonstrates both ordinary and hostile messages:

\codefile{../code/final/tests/test_models_parser.py}

Feeding one byte at a time is more deterministic than trying to force TCP to
produce one-byte reads. Putting two messages in one feed verifies leftover
preservation. Declaring an oversized body without sending it verifies early
rejection. Duplicate lengths and transfer coding test ambiguity policy.

Every accepted feature should have malformed neighbors. If one Content-Length
is accepted, test duplicate, negative, nondecimal, too large, incomplete, and
combined transfer encoding. If percent encoding is accepted, test invalid hex
and invalid UTF-8.

\section{Application and API tests}

Application tests prove static-route precedence, parameter capture, HEAD fallback,
OPTIONS, 404/405, Allow, and middleware nesting without a socket. Task tests use
a fresh temporary SQLite file and call the application directly.

\codefile{../code/final/tests/test_task_api.py}

The CRUD lifecycle test is intentionally stateful within one test: create, read,
patch, replace, delete, and confirm absence. Separate tests cover filters,
pagination, problem structure, query validation, SQL-shaped search text, and
length requirements. Each test case receives its own database, so order does not
matter.

\section{Real socket integration}

\codefile{../code/final/tests/test_server.py}

The test configuration requests port zero. After the readiness event, the test
reads the actual bound address. Teardown calls shutdown and joins the accept
thread even when an assertion fails.

One test uses \code{http.client} twice on a persistent connection and verifies
HEAD body suppression. Another uses raw fragmented writes. Another sends a
syntactically complete request without Host and inspects the 400 problem. The
concurrency test launches 50 clients with explicit deadlines and requires every
one to complete.

\begin{pitfall}
Never let a network test wait forever. Put timeouts on clients, wait for server
readiness, and stop/join the server in teardown. A test that can hang the suite
is itself a reliability defect.
\end{pitfall}

\section{Failure injection checklist}

Before declaring the capstone complete, test:

\begin{itemize}
  \item connection close before header and before declared body completion;
  \item bare LF, folded fields, invalid names, control characters, and missing Host;
  \item unsupported version/method/target form/transfer coding/media type;
  \item limits at one byte below, exactly at, and one byte above each boundary;
  \item two buffered requests and a request split at every byte;
  \item HEAD, 204, OPTIONS, Allow, and Connection close;
  \item invalid JSON encoding/grammar/top-level type/fields/date/enums;
  \item not found, PUT defaults, PATCH emptiness, delete twice;
  \item database constraint and busy failures without partial writes;
  \item queue saturation, read timeout, peer reset, and shutdown during activity.
\end{itemize}

Some OS timing failures can be difficult to make deterministic. Prefer injecting
a fake or using socket pairs when the unit boundary permits. For the few end-to-
end cases, use generous but finite deadlines and assert outcomes rather than
sleeping for guessed durations.

\section{Hardening review}

The implementation has several deliberate defenses: loopback default, strict
CRLF, Host requirement, framing ambiguity rejection, finite bytes/fields/time,
finite workers/queue/connection requests, response-field control validation,
media and UTF-8 checks, allow-list domain validation, parameterized SQL, generic
500 responses, and no toy authentication.

Remaining risks include an incomplete general HTTP implementation, limited
timeout granularity, no TLS, no rate limiting, no authentication/authorization,
no process isolation, minimal signal handling, simple logs, one-process capacity,
and no security maintenance process. Naming these is part of hardening.

\section{One-command verification}

From the repository root:

\begin{wiremessage}
python build.py test
\end{wiremessage}

The helper runs snapshot tests and final tests. \code{python build.py all} also
compiles the book, fails on undefined references and overfull boxes, and copies
the approved PDF into \filepath{dist}.

\begin{checkpoint}
The capstone passes when a clean checkout can run all tests and build the PDF;
50 concurrent local clients finish within deadlines; invalid framing never
reaches a handler; temporary databases leave no persistent test state; and the
reader can explain every intentional unsupported feature.
\end{checkpoint}

\summarytitle

Tests mirror architectural boundaries. Parser tests control exact fragments;
application tests bypass sockets; API tests isolate databases; network tests use
ephemeral ports and deadlines. Hardening is the combination of strict parsing,
finite resources, safe errors, storage discipline, and honest limitations.

\reviewtitle
\begin{enumerate}
  \item Why is byte-at-a-time parser feeding better than hoping TCP fragments that way?
  \item Which properties can only a real socket integration test establish?
  \item Why must tests use temporary database paths and ephemeral ports?
  \item Name five remaining production risks after the suite passes.
\end{enumerate}

\begin{exercise}[Boundary expansion]
Choose one configured limit and add three tests: just below, exactly at, and just
above it. State whether rejection can occur before the complete request arrives.
\end{exercise}

\begin{exercise}[Mini threat model]
List assets, entry points, trust boundaries, attacker capabilities, existing
controls, and missing controls for the Task API. Prioritize five risks by impact
and likelihood without pretending the educational server will be deployed.
\end{exercise}

